\documentclass[hyperref={pdfpagelabels=false}, xelatex, aspectratio=169, 10pt]{beamer}

% 编码设置
\usepackage{fontspec}
\usepackage{ctex} % 中文支持
\usepackage{xeCJK} % 更好的中文支持

% 宏包
\usepackage{amsmath}
\usepackage{amssymb}
\usepackage{amsfonts}
\usepackage{mathrsfs}
\usepackage{mathtools}

% 颜色设置 (可自定义)
\definecolor{MyBlue}{RGB}{43, 87, 151}
\definecolor{MyGray}{RGB}{235, 235, 235}
\definecolor{MyWhite}{RGB}{255, 255, 255}

% 主题设置
\usetheme{Madrid} % 使用 Madrid 主题

% 自定义颜色
\setbeamercolor{palette primary}{bg=MyBlue,fg=MyWhite}
\setbeamercolor{palette secondary}{bg=MyBlue,fg=MyWhite}
\setbeamercolor{palette tertiary}{bg=MyBlue,fg=MyWhite}
\setbeamercolor{palette quaternary}{bg=MyBlue,fg=MyWhite}
\setbeamercolor{structure}{fg=MyBlue} % 结构元素颜色
\setbeamercolor{title in head/foot}{bg=MyBlue,fg=MyWhite} % 标题颜色
\setbeamercolor{section in head/foot}{bg=MyBlue,fg=MyWhite} % 章节颜色
\setbeamercolor{subsection in head/foot}{bg=MyBlue,fg=MyWhite} % 子章节颜色

% 帧标题背景颜色
\setbeamercolor{frametitle}{bg=MyBlue,fg=MyWhite}

% 块(block)颜色
\setbeamercolor{block title}{bg=MyGray,fg=MyBlue}
\setbeamercolor{block body}{bg=MyGray,fg=black}

% 字体设置
\setCJKmainfont{宋体} % 可根据需要替换为其他字体，例如：思源宋体
\setCJKsansfont{黑体} % 可根据需要替换为其他字体，例如：思源黑体
\setCJKmonofont{楷体} % 可根据需要替换为其他字体，例如：思源楷体

% 页脚设置
\setbeamertemplate{footline}
{
  \usebeamerfont{footline}
  \insertframenumber/\inserttotalframenumber
  \hfill%
  \insertshortauthor\ \ \ \ \ \insertshortdate
  \vskip2pt%
}

% 封面页
\title{牛顿法（Newton 法）} % 幻灯片标题
\author{柯力洲} % 作者
\date{2025年9月17日} % 修改为当前日期或您需要的日期

\begin{document}

\frame{\titlepage}

\begin{frame}[fragile]
  \frametitle{牛顿法（Newton 法）}
  \begin{example}[例 8 (Newton 法)]
    设 $f(x)$ 是定义在 $[a, b]$ 上的二次连续可微的
    实值函数，设 $\hat x \in (a, b)$，使得 $f(\hat x) = 0, f'(\hat x) \neq 0$. 求证：存在 $\hat x$ 的邻域
    $U(\hat x)$，使得任选 $x_0 \in U(\hat x)$，迭代序列
    \[
      x_{n+1} = x_n - \frac{f(x_n)}{f'(x_n)} \quad (n=0, 1, 2, \dots)
    \]
    是收敛的，并且
    \[
      \lim_{n \to \infty} x_n = \hat x.
    \]
  \end{example}
  
  \begin{proof}[证]
    考虑不动点迭代函数 $Tx \stackrel{\mathrm{def}}{=} x - \frac{f(x)}{f'(x)}$.
    我们的目标是证明 $T$ 在 $\hat x$ 的某个邻域上是一个压缩映射, 并应用\textbf{巴拿赫不动点定理}。

    \medskip % 增加一点垂直间距

    首先, 计算其导数 $T'(x)$:
  \end{proof}
\end{frame}

\begin{frame}[fragile]
  \frametitle{证明（续）}
    \begin{proof}[续]


    \[
      T'(x) = \frac{d}{dx}(Tx) = 1 - \frac{(f'(x))^2 - f(x)f''(x)}{(f'(x))^2} = \frac{f(x)f''(x)}{(f'(x))^2}.
    \]

    因为 $f(\hat x) = 0, f'(\hat x) \neq 0$, 且 $f, f', f''$ 连续, 所以 $T'(x)$ 在 $\hat x$ 点附近连续, 且 $T'(\hat x) = 0$.

    \medskip

    根据极限的定义, 对于任意给定的 $\alpha \in (0, 1)$, 存在 $\delta > 0$, 使得当 $x \in I = [\hat x - \delta, \hat x + \delta]$ 时, 都有 $|T'(x)| \leq \alpha$.

    \medskip

        我们考虑度量空间 $(I, d)$, 其中 $I = [\hat x - \delta, \hat x + \delta]$ 是实数集 $\mathbb{R}$ 上的闭区间, $d(x, y) = |x - y|$ 是标准的欧氏距离。作为 $\mathbb{R}$ 上的一个闭子集, \textbf{该度量空间是完备的}。

    \medskip

    现在验证压缩映像原理的条件:

        1. \textbf{$T$ 是 $I$ 到自身的映射 ($T: I \to I$)}:
        对任意 $x \in I$, 因为 $T(\hat x)=\hat x$, 由中值定理可得:
        \[ |Tx - \hat x| = |Tx - T\hat x| = |T'(\xi)|\,|x - \hat x| \le \alpha |x - \hat x| \le \alpha\delta < \delta \]
        其中 $\xi$ 在 $x$ 和 $\hat x$ 之间。这说明 $Tx$ 也在 $I$ 内。

        \medskip
        2. \textbf{$T$ 是 $I$ 上的压缩映射}:
        对任意 $x, y \in I$, 由中值定理:

  \end{proof}
\end{frame}

\begin{frame}[fragile]
  \frametitle{证明（续）}
    \begin{proof}[续]

        \[ |Tx - Ty| = |T'(\zeta)|\,|x-y| \le \alpha\,|x-y| \]
        其中 $\zeta$ 在 $x$ 和 $y$ 之间。

    \medskip

    根据\textbf{巴拿赫不动点定理}, 
    
    $T$ 在完备度量空间 $I$ 上存在唯一的不动点 $x^*$。
    
    任取 $x_0 \in I$, 迭代序列 $x_{n+1}=Tx_n$ 必收敛到该不动点。

    该不动点 $x^*$ 满足 $x^* = T(x^*)$, 即 $f(x^*)=0$. 
    
    由于在邻域 $I$ 中 $\hat x$ 是唯一的根, 故 $x^* = \hat x$.

    因此, $\lim\limits_{n \to \infty} x_n = \hat x$.
  \end{proof}
\end{frame}

\end{document}